\chapter{Conclusion and Future Work}
\chaptermark{Conclusion and Future Work}
\HRule\\[2cm]
\label{Chapter6}
\lhead{\emph{\textbf{Chapter 6:} Conclusion and Future Work}}

\begin{spacing}{1.6}

\section{Summary of Work}

This thesis addressed the fundamental conflict between cloud-based anomaly
detection utility and data privacy: the server must see the data to detect
anomalies, yet the data is too sensitive to share. The proposed resolution is
a complete Privacy-Preserving Anomaly Detection framework built on CKKS
Homomorphic Encryption, which enables the detection server to compute anomaly
scores directly on ciphertexts without accessing any plaintext data.

The work spans five major contributions:

\subsection{FHE Library Selection}

A systematic four-experiment benchmarking study of Microsoft SEAL, IBM HElib,
and Pyfhel characterized the performance properties relevant to anomaly
detection workloads: ciphertext multiplication latency, rotation latency,
maximum multiplication depth, and noise budget evolution. The study identified
SEAL as the most performant library for real-valued ML computation, achieving
the deepest multiplication budget (depth 32 at polynomial degree 32768), the
most consistent HEXL hardware acceleration (15--45\% improvement), and the
best arithmetic performance for the ct$\times$pt operations that dominate
anomaly scoring computation.

The benchmarking also established the fundamental asymmetry between addition
and multiplication in HE: addition consumes negligible noise budget (enabling
over 1,048,576 successive additions), while each multiplication consumes
approximately $1/L$ of the remaining budget. This finding directly motivates
the algorithm design principle: every multiplication must be justified, while
rotations and additions are essentially free.

\subsection{Dataset-Agnostic Preprocessing Pipeline}

A five-stage preprocessing pipeline converts raw multivariate data from any
domain to CKKS-encrypted form. The pipeline handles continuous and binary
feature types with type-appropriate CKKS-compatible normalization, enforces
strict separation of normalization parameters from anomaly data, applies
SIMD-optimized column-major and row-major encoding strategies that minimize
HE depth consumption and noise accumulation, and verifies output integrity
through five automated checks. The pipeline was successfully applied to two
benchmark datasets --- SWaT (51 features) and BATADAL (43 features) ---
demonstrating dataset generality.

\subsection{Multi-Round Interactive Protocol for Exact Encrypted Division}

The central technical contribution of this thesis is the multi-round interactive
protocol that resolves the encrypted division problem --- the primary obstacle
in implementing anomaly detection algorithms under CKKS. The protocol routes
the encrypted denominator through the key-holding data owner for exact plaintext
inversion and re-encryption, achieving zero approximation error at zero
additional HE depth cost.

Two alternative approaches --- Newton-Raphson iterative inversion and Chebyshev
polynomial approximation --- were systematically evaluated and shown to be
unviable. Newton-Raphson diverges because the $37\times$ range of feature
variances in real-world data prevents a single initial guess from satisfying
the convergence condition for all SIMD slots simultaneously. Chebyshev
approximation produces unacceptable 33\% error at degree 3 and scale overflow
at higher degrees. The multi-round protocol is the only approach that achieves
exact division, and its privacy guarantee is preserved because only aggregate
model statistics (not individual data points) are revealed to the data owner
during inversion rounds.

\subsection{HE-SAD and HE-PCA Implementations}

Two complete CKKS implementations of unsupervised anomaly detection were
developed and evaluated:

\textbf{HE-SAD} implements variance-normalized Mahalanobis distance scoring
through a 9-round protocol. The per-row computation consumes only 2 of 5
available HE multiplication levels: one for squaring and one for variance
normalization. Six rotations implement the feature slot sum. HE-SAD achieves
AUC $= 0.9533$ on SWaT, matching the plaintext baseline of $0.929$ and
demonstrating that no fundamental accuracy penalty results from privacy-preserving
encryption.

\textbf{HE-PCA} implements PCA reconstruction error scoring through an 8-round
protocol. The per-row computation consumes 3 of 5 levels: one for projection
onto each component, one for reconstruction, and one for the squared residual
norm. Optimal component selection ($k^*$ explaining at least 90\% of training
variance) is performed automatically by the server. A slight AUC reduction
compared to the plaintext PCA baseline is observed, attributable to CKKS
approximation noise accumulating across the $k^*$ component projection
operations.

\subsection{Network Deployment and Live Monitoring}

The complete system is deployed as a three-process network application:

\begin{itemize}
  \item A \textbf{FastAPI detection server} (14 REST endpoints) that receives
  encrypted ciphertexts, performs all HE computation, and computes scores
  asynchronously in a background thread.
  \item A \textbf{data owner client} embedded in the monitoring dashboard,
  handling encryption, inversion rounds, score decryption, and classification.
  \item A \textbf{live SocketIO monitoring dashboard} with browser-based
  controls (Send Data, Pause, Stop, Run SAD, Run PCA, Reset) eliminating the
  need for a separate client terminal.
\end{itemize}

The deployment addresses practical engineering challenges including batched
ciphertext transfer (100 ciphertexts per HTTP request to avoid socket buffer
overflow), asynchronous score computation with client polling, and real-time
monitoring of round progress, per-round latency, score distributions, and
classification metrics.

\section{Key Findings}

\textbf{Finding 1: Exact encrypted division is achievable without approximation
error.}
The multi-round interactive protocol demonstrates that HE-based ML algorithms
requiring division can be implemented exactly by routing aggregate parameter
inversion through the key-holding data owner. This is superior in every
quantified dimension to Newton-Raphson (diverges) and Chebyshev approximation
(33\% error or scale overflow).

\textbf{Finding 2: No fundamental accuracy penalty from encryption for SAD.}
HE-SAD achieves AUC $= 0.9533$ versus plaintext AUC $= 0.929$ on SWaT. The
multi-round protocol's exact division preserves the full discriminative power
of variance-normalized anomaly scoring under CKKS encryption. The slight AUC
improvement over the plaintext baseline is within test set variability and
does not indicate a systematic advantage.

\textbf{Finding 3: CKKS approximation noise accumulation is algorithm-specific.}
HE-SAD shows no accuracy loss while HE-PCA shows a slight reduction. This
difference reflects the number of multiply-plain + rescale operations per row:
SAD requires 2 (one squaring, one normalization), while PCA requires $2k^* + 1$
(two per component plus one for the final squaring). The accumulated $10^{-9}$
approximation error per operation becomes measurable in the score distribution
when dozens of operations are chained.

\textbf{Finding 4: Depth efficiency enables practical leveled HE without
bootstrapping.}
By designing SAD to use 2 levels and PCA to use 3 levels per row, both
algorithms fit comfortably within a 5-level budget at degree 16384, with no
bootstrapping required. This makes the system practical with current SEAL
parameter settings and hardware.

\textbf{Finding 5: PCA's advantage on BATADAL motivates HE-PCA for correlated
attacks.}
On BATADAL, where attacks preserve individual sensor statistics while disrupting
cross-sensor correlations, PCA (F1 $= 0.402$) substantially outperforms SAD
(F1 $= 0.256$). This demonstrates that HE-PCA addresses a qualitatively
different and harder class of anomalies than HE-SAD, and that both algorithms
are needed for a comprehensive privacy-preserving detection framework.

\textbf{Finding 6: Network transport, not HE computation, dominates latency.}
Test data upload accounts for approximately 74\% of total protocol latency
(35 of 47 seconds) due to the large ciphertext size (1 MB each). Pure HE
computation rounds total approximately 6--7 seconds. Cloud deployment over a
dedicated high-bandwidth network would eliminate the transport bottleneck,
making total protocol latency competitive with plaintext detection systems.

\section{Limitations}

\begin{enumerate}
  \item \textbf{HE-PCA accuracy loss.}
  Across the $k^*$ component projections, CKKS approximation noise accumulates
  and causes a measurable AUC reduction compared to the plaintext PCA baseline.
  Bootstrapping (supported by HElib but not the current SEAL configuration)
  would refresh the noise budget to enable deeper circuits and potentially
  recover accuracy.

  \item \textbf{No bootstrapping.}
  The leveled HE configuration supports a fixed number of multiplications per
  ciphertext chain. Algorithms requiring more than 5 levels per row --- such
  as multi-layer neural network inference --- cannot be implemented without
  bootstrapping.

  \item \textbf{Localhost evaluation.}
  All latency measurements were made with client and server on the same machine.
  Real network deployment introduces additional transport latency depending on
  bandwidth, round-trip time, and geographic separation. While the architecture
  is cloud-ready (changing \texttt{SERVER\_HOST} in configuration suffices),
  the performance characterisation is limited to localhost.

  \item \textbf{Binary classification only.}
  The system outputs a binary Normal/Anomalous decision. Multi-class anomaly
  categorization (identifying the type or severity of anomaly) would require
  labeled training data and fundamentally different algorithms.

  \item \textbf{BATADAL HE evaluation pending.}
  The preprocessing pipeline successfully processed BATADAL (43 features,
  8761 normal training rows), and the framework is configured to evaluate it
  with a 3-line configuration change. Full end-to-end HE evaluation on BATADAL
  was not completed within the scope of this thesis.

  \item \textbf{Honest-but-curious adversary model.}
  The system provides strong privacy guarantees against an honest-but-curious
  server. Protection against a malicious server that deviates from the protocol
  (e.g., by returning incorrect ciphertexts) requires additional verifiability
  mechanisms not included in this work.
\end{enumerate}

\section{Future Work}

\begin{enumerate}
  \item \textbf{Reduce HE-PCA accuracy loss.}
  Investigate higher-precision CKKS configurations (larger polynomial modulus
  degree, more primes per level), noise-aware component selection strategies
  (fewer components to reduce accumulated error), or component-wise bootstrapping
  to recover plaintext PCA accuracy in the encrypted setting.

  \item \textbf{Complete BATADAL HE evaluation.}
  Apply both HE-SAD and HE-PCA to BATADAL using the existing framework
  infrastructure. Given PCA's advantage on BATADAL in plaintext, HE-PCA
  evaluation is particularly important to characterise privacy-preserving
  detection of coordinated anomalies.

  \item \textbf{Cloud deployment and latency profiling.}
  Deploy the detection server on cloud infrastructure (AWS EC2 or GCP Compute
  Engine) and measure real network latency for each protocol round across
  different network conditions and geographic separations.

  \item \textbf{ROC-based threshold optimization.}
  Replace the fixed three-sigma threshold rule with ROC-curve-based optimization
  to maximize F1 score or a user-specified precision-recall operating point.

  \item \textbf{Streaming anomaly detection.}
  Extend from batch mode (encrypting all test observations before scoring)
  to streaming mode (scoring each new observation as it arrives). This requires
  incremental update protocols for the mean and variance under encryption.

  \item \textbf{Bootstrapping for deeper circuits.}
  Integrate HElib bootstrapping to support more complex detection algorithms
  --- including shallow neural network autoencoders --- that require more than
  5 HE multiplication levels per inference.

  \item \textbf{Federated privacy-preserving detection.}
  Extend to the federated setting where multiple data owners contribute
  encrypted training data to a shared anomaly detection model without exposing
  their data to each other or to the server. This requires secure aggregation
  protocols for the mean and covariance computation steps.

  \item \textbf{Verifiable computation.}
  Add zero-knowledge proof mechanisms to verify that the server performed the
  correct HE computations, extending the privacy guarantee from honest-but-
  curious to malicious adversaries.

  \item \textbf{Application to additional domains.}
  Evaluate the framework on healthcare (patient physiological monitoring),
  financial (credit card transaction fraud), and network security (intrusion
  detection) datasets to demonstrate that the privacy-preserving detection
  capability generalizes beyond the water infrastructure benchmarks used in
  this thesis.
\end{enumerate}

\section{Conclusion}

This thesis demonstrates that high-quality, unsupervised anomaly detection on
sensitive data is achievable under full homomorphic encryption, with no
fundamental accuracy penalty compared to plaintext detection. The work makes
a concrete contribution at each layer of the privacy-preserving detection
stack: from FHE library selection and dataset preprocessing, through the
central technical innovation of exact encrypted division via the multi-round
interactive protocol, to complete HE implementations of two anomaly detection
algorithms, and finally to a deployed network system with real-time monitoring.

The multi-round interactive protocol is the defining technical contribution.
It demonstrates that the encrypted division problem --- which blocks the
straightforward application of HE to most anomaly detection algorithms ---
can be resolved exactly and efficiently by using the data owner as a trusted
computation agent for aggregate parameter inversion. This insight is general
and applicable beyond the specific algorithms and datasets considered in this
thesis.

HE-SAD achieves AUC $= 0.9533$ on the SWaT benchmark, surpassing the plaintext
baseline of $0.929$. The detection server computed all anomaly scores directly
on CKKS ciphertexts without accessing any plaintext sensor value. The data
owner's secret key never left the client, and no individual measurement was
exposed to the server at any protocol step.

The results establish that privacy-preserving anomaly detection is not merely
theoretically feasible but practically achievable with current homomorphic
encryption technology. The data owner retains complete data sovereignty; the
detection server gains no knowledge of individual measurements; and detection
accuracy is preserved. This resolves the fundamental conflict between anomaly
detection utility and data privacy that motivated this work.

\end{spacing}
